\subsection{研究背景与意义}

近年来，协作机器人与自主移动机器人在工业制造、仓储物流、医疗服务等领域加速落地，
成为推动产业智能化转型的核心载体。
视觉—语言—动作（Vision-Language-Action, VLA）大模型将视觉感知与动作决策统一到端到端神经网络中，
显著拓展了机器人在开放场景下的泛化能力。
国家"十五五"规划（2026—2030年）明确将具身智能机器人列为战略性新兴产业重点方向，
标志着机器人从辅助性工具向智能制造和数字经济核心基础设施的历史性跃升。
在此背景下，如何在资源受限的边缘平台上高效、可靠地运行机器人智能软件，
已成为制约机器人大规模部署的关键问题。

从技术栈角度看，机器人系统通常由三个层次组成。
\textbf{硬件层}涵盖上位机（嵌入式计算板）、传感器（摄像头、力传感器）、执行器（电机、舵机）与机械结构，
决定机器人的物理感知与运动能力；
\textbf{算法层}包括视觉感知、策略规划与运动控制，
以 ACT\cite{act}、Diffusion Policy\cite{diffusionpolicy} 等端到端模仿学习模型为代表；
\textbf{基础软件层}则由操作系统（Linux 等）、通信中间件（ROS 等）和推理框架（PyTorch、ONNX Runtime、TFLite 等）构成，
向下管理硬件资源，向上支撑算法模型的运行。
本文的研究聚焦于基础软件层：操作系统的内核调度策略、系统调用路径与设备驱动模型，
以及推理框架的线程管理与算子调度，
共同决定着机器人系统的硬件生态兼容性、推理性能、决策准确性、运行可靠性和系统安全性，
是整个机器人软件栈性能与可信度的基础。

面向机器人基础软件的研究已有大量探索，但仍存在明显不足。
广义上，机器人基础软件涵盖三个方向：操作系统内核、推理库、以及专用于机器人的通信工具。

在\textbf{操作系统内核}层面，现有路线可归纳为三类：
以 Linux 为代表的\textbf{宏内核系统}应用广泛，但系统调用需经历用户态—内核态切换，
单次操作引入百纳秒至数微秒的不确定延迟；
以 HarmonyOS、Fuchsia 为代表的\textbf{微内核系统}结构精简、隔离性强，
但驱动驻于用户态，硬件操作需经历两次态切换，整体延迟往往反而高于宏内核；
以 FreeRTOS、Zephyr 为代表的\textbf{极简嵌入式实时系统}调度开销极低，
但缺乏完整的深度学习推理支持，难以直接承载现代策略模型。

在\textbf{推理库}层面，PyTorch、ONNX Runtime、TFLite 各有侧重，
均以"外挂式"方式接入操作系统：
其内置线程池（OpenMP 或 C10）的调度与 OS 内核 CFS 调度器相互独立，
在资源受限平台上容易引发调度冲突，造成额外的上下文切换与 futex 等待开销。

在\textbf{机器人专用通信工具}层面，ROS/ROS2 以 DDS 为底层通信层，
micro-ROS 则将其下沉至嵌入式 MCU，二者均简化了多节点协作开发，
但本质上继承了底层内核的调度局限，节点间消息传递延迟受制于 Linux CFS 调度器，无硬实时保证。

更根本的问题在于，现有性能优化与分析工作多聚焦于单一层次。
主流优化路径往往集中在模型压缩、模型结构改进或硬件通信链路调优上，
而对推理运行时、Python 线程行为、OS 内核调度、系统调用阻塞与设备 I/O 栈等
基础软件因素关注不足。
这使得许多工作虽然能够报告端到端延迟或单点优化收益，
却缺乏从 AI 推理框架层到 OS 内核层的跨层次联合剖析，
难以回答"瓶颈究竟在哪一层、由什么原因导致"等关键问题。

针对上述研究缺口，本文以 HuggingFace LeRobot 框架\cite{lerobot}下
ACT 策略在树莓派5（ARM Cortex-A76，4核，4\,GB LPDDR4X，无 GPU）上的在线推理为研究对象，
系统开展机器人基础软件的跨层次性能剖析。
实验表明，单次 ACT 完整推理耗时 1924.8\,ms，
其中 ResNet18 视觉编码器占 66.6\%（1281\,ms）为第一瓶颈；
在 chunk\_size\,=\,100 的同步推理配置下实测帧率仅 18.7\,fps，距 30\,fps 目标仍有约 38\% 缺口。
通过推理异步化优化（hold=0.30），有效帧率提升至 27.42\,fps（相对提升 70\%），基本接近实时目标。
本文提出覆盖 AI 推理框架层、Python 运行时层、OS 内核层和硬件 I/O 层的四层性能分析模型，
建立面向"感知—推理—执行"紧耦合场景的多工具联合剖析方法，
通过 torch.profiler/ftrace/perf 多工具联合剖析，端到端量化各阶段时间开销，
定位 ResNet18 编码器为第一瓶颈（占 66.6\%），
揭示 strace 带来 40.9\% FPS 下降，验证 ftrace 的低扰动必要性，
为嵌入式机器人基础软件的性能优化提供明确的优先级依据。
